\documentclass[11pt,a4paper]{article}
\usepackage[margin=1in]{geometry}
\usepackage{booktabs}
\usepackage{longtable}
\usepackage{hyperref}
\usepackage{xcolor}
\usepackage{soul}
\usepackage{enumitem}

\title{Replication Report: Ghattargi et al. (2018) \\
\large Comparative genome analysis reveals key genetic factors associated with probiotic property in \emph{Enterococcus faecium} strains}
\author{Ollie (OpenClaw AI) --- BVBRC Replication Project, Wave 4 target \#16}
\date{2026-06-25 (promoted from spot-check 2026-06-17)}

\begin{document}
\maketitle

\begin{abstract}
We independently re-tested the central safety and genome-stability claims of Ghattargi \emph{et al.} (\emph{BMC Genomics} 19:652, 2018; \href{https://doi.org/10.1186/s12864-018-5043-9}{doi:10.1186/s12864-018-5043-9}) for the candidate probiotic \emph{E. faecium} strain 17OM39 (BV-BRC \texttt{1352.1047}) against the marketed probiotic T110 (\texttt{1344042.3}). Using BV-BRC's specialty-gene tables (CARD, NDARO, VFDB, Victors, PATRIC) and full-CDS product scans over 5{,}776 (17OM39) and 5{,}173 (T110) annotated genes, we \textbf{reproduce} the paper's core AMR-absence claim (no \emph{vanA}/\emph{vanB}, no canonical \emph{tet}) and \textbf{substantially reproduce} the virulence-absence claim (17OM39 is in fact \emph{cleaner} than T110, which carries the \emph{Cna} collagen adhesin). We \textbf{do not reproduce} the mobile-genetic-element (MGE) stability claim: BV-BRC RASTtk annotation shows 17OM39 with $\sim 2.7\times$ more MGE-like CDS than T110, opposite the paper --- but this comparison is confounded by 17OM39 being a 106-contig draft vs.\ T110 being finished. The core-genome phylogeny (C6) was not re-run. \textbf{Verdict: PARTIAL} --- 4 of 6 testable claims reproduced, 1 explicit miss with a known confound, 1 untested.
\end{abstract}

\section{Paper Under Test}
Ghattargi \emph{et al.} compare the isolate 17OM39 (draft; GCF\_001652715.1; 2{,}840{,}201 bp; 106 contigs) against a marketed probiotic T110 (finished; GCA\_000737555.1; 2{,}737{,}963 bp + 44\,kb plasmid), a non-pathogenic non-probiotic (NPNP) strain, and pathogenic comparators. The paper concludes 17OM39 is a credible probiotic candidate because it (i) clusters with T110 on the core-genome tree, (ii) lacks \emph{vanA}/\emph{vanB} and tetracycline resistance genes, (iii) lacks functional virulence genes, and (iv) has higher genomic stability (fewer transposons/MGEs).

\section{Claims Tested}
\begin{longtable}{cll}
\toprule
\# & Claim & Tested this pass? \\
\midrule
C1 & 17OM39 + T110 deposited and adequately sequenced & \checkmark (metadata) \\
C2 & 17OM39 genome size / topology consistent with \emph{E.\ faecium} & \checkmark (metadata) \\
C3 & 17OM39 lacks vancomycin / tetracycline resistance genes & \checkmark BV-BRC sp\_gene + product scan \\
C4 & 17OM39 lacks functional virulence genes & \checkmark BV-BRC sp\_gene (VFDB/Victors) \\
C5 & 17OM39 has fewer MGEs than pathogenic strains & partial --- 17OM39 vs T110 only, \textbf{opposite direction} \\
C6 & 17OM39 clusters with T110 on core-genome phylogeny & \textbf{not run} (needs Roary/PhyloPhlAn) \\
\bottomrule
\end{longtable}

\section{Method}
\subsection{BV-BRC specialty-gene rescreen}
For both genomes we pulled the full BV-BRC \texttt{sp\_gene} table:
\begin{verbatim}
curl 'https://www.bv-brc.org/api/sp_gene/?eq(genome_id,<GID>)
      &select(genome_id,property,source,gene,product,evidence,classification)
      &limit(2000)'
\end{verbatim}
These are the public, pre-computed AMR + virulence calls BV-BRC maintains by aligning each genome's annotated proteins against CARD, NDARO (NCBI AMRFinder), VFDB, Victors, and PATRIC's own AMR/VF curations --- the same reference DBs the paper used (CARD/ResFinder for AMR, VFDB for virulence).

\subsection{Full-CDS product scan}
We also pulled every CDS feature (\texttt{/genome\_feature/?...feature\_type=CDS}) and scanned product names locally with regex for MGE, AMR and probiotic/virulence keywords. Totals: 5{,}776 CDS (17OM39) and 5{,}173 CDS (T110).

\subsection{What we did NOT run}
Roary/PGAP pan-genome (so no C6 phylogeny); \emph{de novo} ResFinder/CARD-RGI/VFDB BLAST against FASTA; ISEScan/ICEberg \emph{de novo} MGE annotation.

\section{Results}
\subsection{Genome metadata (C1, C2)}
Both deposited, both sized correctly for \emph{E. faecium}. 17OM39: 2{,}840{,}201 bp / 106 contigs (draft). T110: 2{,}737{,}963 bp finished + 44\,kb plasmid \texttt{1344042.14}. \checkmark

\subsection{AMR rescreen (C3) --- REPRODUCED}
BV-BRC \texttt{sp\_gene}, property = ``Antibiotic Resistance'':

\begin{longtable}{lrrrrrrr}
\toprule
Genome & Total sp\_gene & AMR rows & CARD & NDARO & PATRIC & \texttt{van*} & \texttt{tet*} \\
\midrule
17OM39 & 158 & 56 & 18 & 3 & 35 & \textbf{0} & \textbf{0} \\
T110   & 148 & 52 & 14 & 3 & 35 & \textbf{0} & \textbf{0} \\
\bottomrule
\end{longtable}

Local product-name scan across all CDS:
\begin{itemize}[nosep]
\item \textbf{vancomycin / vanA / vanB / vanC}: 0 hits in either genome. \checkmark
\item \textbf{tetracycline / tetM / tetO / tetL}: 17OM39 has \emph{one} CDS annotated ``tetracycline resistance MFS efflux pump'' (no canonical gene symbol, not a ribosomal-protection \emph{tet} gene); T110 has zero. Minor discrepancy with the paper's ResFinder run.
\end{itemize}

The AMR rows that \emph{are} present in both genomes are nearly all intrinsic/housekeeping alignments (EF-Tu, EF-G, gyrase B, ileS, DHFR, D-Ala-D-Ala ligase, LiaFSR). The only ``acquired-style'' AMR calls shared by both --- \texttt{Lsa(A)} and \texttt{Msr(C)} ABC-F ribosomal protection --- are known to be \textbf{chromosomally intrinsic to \emph{E. faecium}}, fully consistent with the paper's ``no acquired AMR'' conclusion. \textbf{C3 reproduces.}

\subsection{Virulence rescreen (C4) --- MOSTLY REPRODUCED, with a real difference}
\begin{longtable}{lrp{9.5cm}}
\toprule
Genome & VF rows & Notable hits \\
\midrule
17OM39 & 18 & Housekeeping/probiotic-associated only: MalR, ClpP, EF-LepA, MAP, thymidylate synthase, \textbf{choloylglycine hydrolase (= bile salt hydrolase, a probiotic trait)}, PerR. \emph{No collagen-binding protein, no Esp, no cytolysin, no gelatinase, no aggregation substance.} \\
T110   & 19 & Same housekeeping set \textbf{plus Collagen-binding protein Cna (VFDB)} --- a \emph{bona fide} \emph{E.\ faecium} virulence/adhesion factor. \\
\bottomrule
\end{longtable}
The paper's claim that 17OM39 lacks functional virulence genes reproduces. In fact 17OM39 is \emph{cleaner} on this axis than T110 --- a small positive safety finding the paper did not emphasize.

\subsection{MGE comparison (C5) --- NOT REPRODUCED (opposite direction, with caveat)}
\begin{longtable}{lrr}
\toprule
MGE keyword & 17OM39 (5{,}776 CDS) & T110 (5{,}173 CDS) \\
\midrule
transposase & \textbf{66} & 21 \\
mobile element & \textbf{56} & 25 \\
phage & \textbf{47} & 19 \\
integrase & \textbf{26} & 4 \\
recombinase & 18 & 12 \\
resolvase & 10 & 3 \\
site-specific recombinase & 8 & 0 \\
invertase & 3 & 1 \\
prophage & 2 & 0 \\
insertion sequence & 0 & 1 \\
conjugal (TraE) & 1 & 1 \\
\midrule
\textbf{TOTAL MGE-like CDS} & \textbf{237 (4.10\%)} & \textbf{87 (1.68\%)} \\
\bottomrule
\end{longtable}
17OM39 has $\sim 2.7\times$ as many MGE-like CDS as T110, opposite the paper.

\textbf{Assembly-status confound.} 17OM39 is a 106-contig SPAdes draft; T110 is finished. Draft assemblies massively inflate transposase counts because repetitive IS elements fragment across contig boundaries (a single IS is reported as multiple partial CDS). The paper's actual C5 comparison was 17OM39 vs \emph{pathogenic clinical} isolates, which we did not pull. C5 cannot be rejected from this evidence, but the public BV-BRC RASTtk annotation does not, on its face, support the paper's stability claim.

\subsection{Probiotic-associated genes}
Both strains carry a comparable probiotic toolkit (3/3 bile salt hydrolase, 9 vs 13 bacteriocins, 13 vs 9 sortase A, 6/6 LPxTG anchored proteins), consistent with the paper.

\section{Verdict}
\begin{longtable}{cl}
\toprule
Claim & Reproduced? \\
\midrule
C1 (data deposits) & \checkmark Exact \\
C2 (genome size) & \checkmark Exact \\
C3 (no \emph{vanA}/\emph{vanB}, no \emph{tet}) & \checkmark Reproduced (one minor MFS-efflux caveat) \\
C4 (no functional virulence) & \checkmark Reproduced --- 17OM39 cleaner than T110 (no Cna) \\
C5 (fewer MGEs) & \textbf{$\times$} Public BV-BRC RASTtk shows opposite; confounded by draft vs.\ finished \\
C6 (core-genome cluster with T110) & \textbf{not tested} (needs Roary/PhyloPhlAn) \\
\bottomrule
\end{longtable}

\textbf{PARTIAL.} 4 of 6 testable claims reproduced; 1 explicit miss with a known confound; 1 untested. Coverage 6/10, Agreement 8/10.

\section{GENUINE CRITIQUE}
This section is deliberately harsh where the evidence warrants; it is the mandatory reproducibility-blocker analysis.

\subsection{The MGE-stability claim (C5) is the paper's weakest link}
The paper does not publish (a) the exact ISEScan/IScompass/equivalent command line and version used to count ``transposons'', (b) the list of pathogenic comparator accessions used as the MGE baseline, or (c) how draft-vs-complete assembly status was controlled. Our re-analysis of public BV-BRC RASTtk finds 17OM39 carrying $\sim 2.7\times$ \emph{more} MGE-like CDS than T110 --- pointing the opposite direction from the paper. Without the paper's exact MGE-counting command this cannot be cleanly adjudicated, and that is the canonical ``underspecified methods'' reproducibility failure.

\subsection{Sample size is two}
A serious probiotic-safety case would compare 17OM39 against $>10$ \emph{E.\ faecium} clinical isolates and $>10$ commercial probiotic isolates. The paper does not, and neither does this pass. Any strong conclusion about ``17OM39 is safer than pathogens'' from an $n=2$ head-to-head is fragile.

\subsection{Annotation-pipeline coupling}
BV-BRC's AMR + VF calls use CARD/NDARO/VFDB but via RASTtk's alignment/threshold choices, not the paper's original ResFinder/CARD-RGI runs. So our agreement tests ``do the same conclusions hold under a \emph{current} pipeline?'' rather than ``do the \emph{exact} numbers reproduce?''. For a safety-relevant claim like C3 this is actually a \emph{stiffer} test, which is why the reproduction is meaningful.

\subsection{C6 phylogeny is the central paper-specific claim and is untested here}
``17OM39 clusters with T110 on the core-genome tree'' is what makes 17OM39 look probiotic-like in the first place. This is what a paper-specific replication should nail down; we punt to a future Roary/PhyloPhlAn pass. Until then, we cannot say the phylogenetic backbone of the paper is reproduced --- only that its safety-relevant leaf claims are.

\subsection{What the paper genuinely got right}
C3 and C4 --- the central safety claims --- are robust under independent BV-BRC reannotation. That is the clinically important part of the paper and it survives. Even better, one \emph{additional} finding surfaces that the paper did not emphasize: T110 carries the \emph{Cna} collagen-binding adhesin (a known \emph{E.\ faecium} virulence factor) and 17OM39 does not, which strengthens rather than weakens 17OM39's safety story.

\subsection{Bottom line}
This paper's clinically load-bearing claims (no acquired AMR, no functional virulence in 17OM39) reproduce. Its genome-stability story (C5) is underspecified, cannot be reproduced from public data as stated, and may not survive an ISEScan rerun on finished-vs-finished comparators. Its phylogenetic story (C6) has not yet been re-tested. PARTIAL is the honest verdict.

\section{Path to Full REPLICATED}
Roughly 10 CPU-hours, not gated by data availability, gated only by having the tools installed:
\begin{enumerate}[nosep]
\item \texttt{prokka} + \texttt{roary -p 8 -i 90 -e --mafft} on 4 paper strains (17OM39, T110, Aus0004 as a pathogenic comparator, one NPNP); build core-gene tree with FastTree; directly tests C6. ($\sim$6--8 CPU-h.)
\item \texttt{ISEScan} on the four FASTAs (not RAST product names) for a fair MGE count that ignores contig-break duplicates; directly re-tests C5. ($\sim$2 CPU-h.)
\item \texttt{abricate --db resfinder} and \texttt{abricate --db vfdb} at 2026 DB versions, confirming C3/C4 with the exact DBs the paper used. ($\sim$10 CPU-min.)
\end{enumerate}

\end{document}
